\section{Conclusion}\label{sec:conclusion}

We tested three biologically-grounded agent reliability features
against naive non-biological alternatives, with 10M+ data points
across 10 independent seeds plus a follow-up experiment with real
sentence embeddings.  All three features demonstrate empirical
benefits under the right conditions.

Metabolic priority gating provides a structural guarantee that
critical operations are served under resource pressure (100\% vs
39.8\% for a flat counter under bursty load).  Autoinducer-based
quorum sensing achieves zero false-positive rate with meaningful
true-positive rate (71--87\% at 40\% compromise), occupying a
precision--recall position unreachable by majority voting or
independent detection.  Bayesian two-signal stagnation detection
dominates on convergence discrimination with real embeddings
(96\% vs 2--40\% naive) but loses with mock embeddings, demonstrating
that biological designs relying on information quality succeed only
when that quality is present.

The consistent pattern: biological design earns complexity through
\emph{mechanism-level structural guarantees}---hard properties enforced
by design (priority gating, signal decay, two-signal discrimination)---rather
than through algorithmic sophistication alone.  Embedding quality
determines whether the two-signal structure delivers its guarantee,
just as metabolic state determines whether priority gating activates.

\paragraph{Open problems.}
Confirmation of the real-embedding result across additional models
and larger sample sizes.  Auto-calibration of quorum sensing thresholds
from population and signal statistics.  Empirical validation of
epistemic topology bounds against measured performance.  End-to-end
evaluation with actual LLM agents performing real tasks.
